\section{Stacks}\label{sec:probstacks}

In this section the stack of Listing~\ref{lst:stack} holds \cc{int}s; characters are
pushed as their codes. ``Stack operations only'' means \cc{stInit}, \cc{stPush},
\cc{stPop}, \cc{stTop}, \cc{stEmpty}, \cc{stSize}: no access to \cc{items[]}.

\problem{32}{Valid brackets}{Easy}
Given a string \cc{s} containing only the characters \cc{( ) [ ] \{ \}}, return 1 if
every bracket is closed by the same kind of bracket and in the correct order, else 0.
\begin{lstlisting}[style=starter]
int isBalanced(char s[]);
\end{lstlisting}
\permitted Stack operations only.
\constraints $0\le|s|\le10^4$.
\required $O(n)$ time.
\begin{tests}
\{[()]\} & 1\\
([)] & 0\\
(( & 0\\
)) & 0\\
(empty string) & 1\\
()[]\{\} & 1\\
\end{tests}

\problem{33}{Infix to postfix and postfix evaluation}{Medium}
(a)~Given an infix expression \cc{in} made of single-digit operands, the operators
\cc{+ - * /} and parentheses (no spaces), write its postfix form into \cc{out}. The
operators \cc{*} and \cc{/} have higher precedence than \cc{+} and \cc{-}; operators of
equal precedence are left-associative.
(b)~Given a valid postfix expression \cc{p} of single digits and \cc{+ - * /}, return its
value (integer division, truncating toward zero).
\begin{lstlisting}[style=starter]
void toPostfix(char in[], char out[]);
int evalPostfix(char p[]);
\end{lstlisting}
\permitted Stack operations only.
\constraints $1\le|\mathit{in}|\le100$; the input is valid.
\required $O(n)$ time for both.
\bonus Explain what changes if \cc{\^} (exponent, right-associative) is added.
\begin{tests}
2+3*4 & 234*+ ; value 14\\
(2+3)*4 & 23+4* ; value 20\\
2*(3+4)-5 & 234+*5- ; value 9\\
8-3-2 & 83-2- ; value 3\\
9-8/(4-2) & 9842-/- ; value 5\\
\end{tests}

\problem{34}{A stack with constant-time minimum}{Easy}
Design a stack that supports \cc{push}, \cc{pop} and \cc{min} (the smallest item
currently on the stack), each in $O(1)$ time.
\begin{lstlisting}[style=starter]
typedef struct { stack all; stack mins; } minStack;
void msInit(minStack *m);
void msPush(minStack *m, int x);
int msPop(minStack *m, int *x);    // 0 if empty
int msMin(minStack *m, int *x);    // 0 if empty
\end{lstlisting}
\permitted Stack operations on the two member stacks.
\required Every operation $O(1)$.
\bonus Use less memory when many pushes are not new minima (push to \cc{mins} only when
$x\le$ the current minimum), and explain why the test must be $\le$, not $<$.
\begin{tests}
push 5 3 7 3 1; min & 1\\
pop; pop; min & 3 (stack 5 3 7)\\
push 2; min; pop; min & 2; 3\\
min on an empty stack & fails (0)\\
\end{tests}

\problem{35}{Next greater element}{Medium}
For every item \cc{a[i]} of an array, find the first item to its right that is
\emph{strictly} greater, or $-1$ if there is none.
\begin{lstlisting}[style=starter]
void nextGreater(int a[], int n, int out[]);
\end{lstlisting}
\permitted Array access to \cc{a} and \cc{out}; one stack.
\constraints $0\le n\le10^5$.
\required $O(n)$ time ($O(n^2)$ earns partial credit).
\bonus The \emph{circular} version: the search may wrap around to the beginning of the
array (\cc{[1,2,1]} $\to$ \cc{[2,-1,2]}).
\begin{tests}
{}[2,1,5,3] & [5,5,-1,-1]\\
{}[2,7,3,5,4,6,8] & [7,8,5,6,6,8,-1]\\
{}[5,4,3] & [-1,-1,-1]\\
{}[1,1] & [-1,-1]\\
\end{tests}

\problem{36}{Daily temperatures}{Medium}
Given the daily temperatures \cc{t[0..n-1]}, compute \cc{wait[i]}, the number of days
after day $i$ until a \emph{strictly} warmer day, or 0 if no warmer day follows.
\begin{lstlisting}[style=starter]
void dailyTemperatures(int t[], int n, int wait[]);
\end{lstlisting}
\permitted Array access to \cc{t} and \cc{wait}; one stack.
\constraints $1\le n\le10^5$; $30\le t[i]\le100$.
\required $O(n)$ time.
\begin{tests}
{}[73,74,75,71,69,72,76,73] & [1,1,4,2,1,1,0,0]\\
{}[30,40,50,60] & [1,1,1,0]\\
{}[30,60,90] & [1,1,0]\\
{}[50,50] & [0,0]\\
\end{tests}

\problem{37}{Online stock span}{Medium}
Prices of a stock arrive one per day. The \emph{span} of today's price is the number of
consecutive days, ending today and going backwards, on which the price was $\le$ today's
price. Implement \cc{ssNext}, which receives today's price and returns its span. The
prices are given one at a time: you cannot look at future prices, and you should not
store the whole history.
\begin{lstlisting}[style=starter]
typedef struct { stack prices; stack spans; } stockSpanner;
void ssInit(stockSpanner *s);
int ssNext(stockSpanner *s, int price);
\end{lstlisting}
\permitted Stack operations on the two member stacks.
\required $O(1)$ amortised per call.
\begin{tests}
100 80 60 70 60 75 85 & 1 1 1 2 1 4 6\\
10 10 10 & 1 2 3\\
5 4 3 2 1 & 1 1 1 1 1\\
\end{tests}

\problem{38}{Asteroid collision}{Medium}
Each item of \cc{a} is an asteroid in a row: its absolute value is its size, its sign its
direction ($>0$ moves right, $<0$ moves left); all move at the same speed. When two
meet, the smaller one explodes; if both are the same size, both explode. Two asteroids
moving in the same direction never meet. Write the survivors, left to right, into
\cc{out} and return how many there are.
\begin{lstlisting}[style=starter]
int asteroidCollision(int a[], int n, int out[]);
\end{lstlisting}
\permitted Array access to \cc{a} and \cc{out}; one stack.
\constraints $0\le n\le10^4$; $a[i]\ne0$.
\required $O(n)$ time.
\begin{tests}
{}[5,10,-5] & [5,10]\\
{}[8,-8] & []\\
{}[10,2,-5] & [10]\\
{}[-2,-1,1,2] & [-2,-1,1,2]\\
\end{tests}

\problem{39}{Decode string}{Medium}
A string is encoded by the rule \cc{k[part]}: the part inside the brackets is repeated
exactly $k$ times ($k\ge1$, possibly several digits). Encodings may be nested. Write the
decoded string into \cc{out}.
\begin{lstlisting}[style=starter]
void decodeString(char s[], char out[]);
\end{lstlisting}
\permitted Character access to \cc{s} and \cc{out}; stacks.
\constraints The input is valid; the decoded string has at most $10^4$ characters.
\required Time proportional to the length of the output.
\begin{tests}
3[a]2[bc] & aaabcbc\\
3[a2[c]] & accaccacc\\
2[abc]3[cd]ef & abcabccdcdcdef\\
10[a] & aaaaaaaaaa\\
abc & abc\\
\end{tests}

\problem{40}{Remove adjacent duplicates}{Easy}
Repeatedly remove two adjacent equal letters from \cc{s} until no two adjacent letters
are equal; write the result into \cc{out}. The result is unique.
\begin{lstlisting}[style=starter]
void removeAdjacent(char s[], char out[]);
\end{lstlisting}
\permitted Character access; one stack.
\constraints $1\le|s|\le10^5$.
\required $O(n)$ time. (Rescanning the string after each removal is $O(n^2)$.)
\begin{tests}
abbaca & ca\\
azxxzy & ay\\
aa & (empty string)\\
abc & abc\\
\end{tests}

\problem{41}{Sort a stack}{Medium}
Sort the stack \cc{s} so that the smallest item is on top. You may use one additional
stack and scalar variables, nothing else.
\begin{lstlisting}[style=starter]
void sortStack(stack *s);
\end{lstlisting}
\permitted Stack operations only; one extra stack.
\constraints $0\le n\le100$.
\required $O(n^2)$ time.
\begin{tests}
bottom..top 3 1 4 1 5 & popping gives 1 1 3 4 5\\
bottom..top 1 2 3 & popping gives 1 2 3\\
empty & empty\\
\end{tests}

\problem{42}{Validate stack sequences}{Medium}
The distinct numbers \cc{pushed[0..n-1]} are pushed onto an initially empty stack in this
order, and pops may occur at any moment. Return 1 if the pops can produce exactly the
order \cc{popped[0..n-1]}, else 0.
\begin{lstlisting}[style=starter]
int validateStackSequences(int pushed[], int popped[], int n);
\end{lstlisting}
\permitted Array reads; one stack.
\constraints $0\le n\le1000$.
\required $O(n)$ time.
\bonus Prove that the greedy rule ``pop whenever the top is the next wanted item'' never
misses a solution.
\begin{tests}
pushed 1 2 3 4 5, popped 4 5 3 2 1 & 1\\
pushed 1 2 3 4 5, popped 4 3 5 1 2 & 0\\
pushed 2 1 0, popped 1 2 0 & 1\\
\end{tests}
